\ofsubsection{Equipmento}
%
\ofquote{"Ah, é mesmo, você não sabia? Hoje em dia, tudo que precisa para realizar seus sonhos é dinheiro e poder."\\}{President Shinra}\\
%
\includegraphics[width=\columnwidth]{./art/images/ff1.jpg}
%
\vfill
%
A força de combate de um personagem pode ser melhorada através de seu \accf{Equipamento}.
Enquanto \accf{Armas} aumentam o dano causado, \accf{Armaduras} protegem contra possíveis danos e \accf{Acessórios} servem de complemento.
Cada equipamento fornece um efeito \accf{Único} e, armas e armaduras têm um \accf{Tipo} adicional de efeito determinado pelo seu tipo.
Armas tem alcance de 1u, exceto se especificado o contrário.
Ao invés de armaduras, os personagens também podem usar roupas comuns que fornecem DEF~1.
Acessórios podem ser usados por todos, mas os personagens só podem equipar armas e armaduras determinadas pela sua profissão.
Embora, os personagens possam usar \accf{Itens}, que forneçam benefícios rápidos dentro e fora do combate, eles são consumidos após seu uso. 
Equipamentos são categorizados em diferentes \accf{Níveis de Equipamento} a depender de sua potência e raridade.
Enquanto todos carregarem equipamento Iniciante, eles precisam alcançar Níveis mais altos para usar equipamentos Avançados e Especialistas.
Os detalhes de cada Nível de Equipamento está disponível na tabela abaixo.
%
\vfill
%
\oftable{p{0.2\columnwidth} p{0.125\columnwidth} p{0.15\columnwidth} p{0.15\columnwidth} p{0.3\columnwidth}}
{\accf{Nível de\newline Equip.} & \accf{Nível} & \accf{Dano da Arma} & \accf{Armadura\newline DEF/RES} & \accf{Valor\newline Aprox.}}
{
	Iniciante & 	1 - 3 	& 1d & +0 / +0   & 500G\ofrow
	Avançado &  4 - 7   & 2d & +1 / +1   & 2000G\ofrow
	Especialista   &  8 - 10	& 3d & +2 / +2   & 4000G 
}
%
\vfill
%
Todos os itens e posses não equipados são armazenadas no \accf{Inventário} do personagem.
Equipamentos e itens podem ser saqueados de adversários derrotados ou recebidos como recompensas por completar missões.
Mas o grupo também pode comprar ou vender produtos específicos em lojas e a mercadores.
A moeda usada para comércio é chamada de \accf{Gil} (\accf{G}) e normalmente vem em forma de moedas de ouro.
Você pode tentar comprar e vender quase qualquer coisa, à medida em que haja alguém disposto a negociar.
%
\newpage
%
\ofboxwithtitle{Exemplo: Comércio}
{
	Terra e seu grupo visitam uma casa de leilão para conseguir itens raros. 
	Hoje é o dia de sorte deles, o item em leilão é uma Chocobo falante!
	O grupo está impressionado com a talentosa criatura e decide usar todas as suas economias, um total de 10.000G!
	Parece que o lance deles está superando ao de todos, mas no último segundo um pai amável anuncia 500.000G e adquire a Chocobo como um presente para seu filho.
	Talvez o grupo tenha mais sorte (ou dinheiro) da próxima vez...
}
%
\vfill
%
\ofquote{"Cloud, assine isso. É um contrato que diz quando a guerra acabar, todas as matérias pertencerão a mim."\\}{Yuffie}
%
\vfill
%
Armas e armaduras podem ser melhoradas através do uso de \accf{Matéria}.
Matéria normalmente parece com um cristal, que concentra forças magicas e pode ser usada por qualquer um.
Elas podem ser encaixadas em espaços nas armas ou armaduras e garantir ao usuário efeitos adicionais.
Aquelas encaixadas também podem ser removidas, trocadas ou armazenadas com facilidade.
A menos que especificado o contrário, cada arma e armadura pode carregar uma matéria, enquanto roupas e acessórios nenhuma.
Matérias podem ser divididas em duas categorias: Comando e Suporte.
Matérias de Comando permitem ao usuário usar novas habilidades, enquanto as de Suporte garantem efeitos passivos.
Se uma Matéria permite o uso de uma habilidade existente, o usuário deve atender às suas exigências, como custo em PM e tempo de conjuração.
%
\vfill
%
\ofquote{"Sim, as doze armas lendárias. São armas. São lendárias. Há até doze delas."\\}{Ghido}
%
\vfill
%
A menos que o MJ diga o contrário, cada personagem tem 1.500G no Nível~1, com esse valor ele pode comprar seus equipamentos iniciais, desta forma:
\ofrow
\ofbullet{Comprar qualquer arma Iniciante por 500G ou uma arma de Mithril por 250G.}
\ofbullet{Comprar uma armadura Iniciante por 500G, uma armadura de Mithril por 250G ou roupas por 100G.}
\ofbullet{Opcional, comprar qualquer conjunto de itens que puder pagar.}
\ofbullet{Opcional, compre um acessório Iniciante por 500G.}
\ofbullet{Opcional, compre uma Matéria que possa pagar.}\vspace*{-0.25cm}\\
As próximas páginas dão vários exemplo de equipamentos para se escolher.
As tabelas detalham efeitos únicos de cada um, assim como os tipos diferentes de armas e armaduras.
Todas as armas e armaduras listadas são do nível Iniciante, a menos que o MJ diga o contrário.
Perceba que a ficha de personagem fornece vários campos para acompanhar o inventário assim como os equipamentos atuais e seus efeitos.
%
\clearpage
%
\newcommand{\ofequipnote}[1]{\multicolumn{2}{p{\columnwidth}}{\hspace*{-0.15cm}\parbox{0.97\columnwidth}{\accf{Efeito Tipo:} #1}}}
\oftable{p{0.22\columnwidth} p{0.7\columnwidth}}
{\oficonweapon\accf{Arcos} & \accf{Efeito Único}}
{
	Arco Mythril & Possui um espaço de Matéria extra. \ofrow
	Arco Élfico & Adiciona 3 ao dano ao atingir um alvo que tenha algum Efeito de Estado.\ofrow
	Arco\newline Assassino & Num Acerto Crítico, o alvo sofre KO. \ofrow
	Besta & Não sofre da penalidade do seu tipo de efeito, mas só tem alcance de 3u. \ofrow
	Aplu & Ignora efeitos que permitem ao alvo ter Vantagem no teste de evasão, exceto seu próprio tipo de efeito. \ofrow
	Anteros & Quando o alvo rolar 5 ou menos no teste de evasão, ele fica Cego por 3 rodadas.\ofrow
	\ofequipnote{Arcos tem alcance de 5u. Se você atacar ou usar uma técnica que precise atacar e se mover no mesmo turno, o alvo tem vantagem no teste de evasão.\phantom{y}}
}
%
\vfill
%
\oftable{p{0.22\columnwidth} p{0.7\columnwidth}}
{\oficonweapon\accf{Arma\newline de Fogo} & \accf{Efeito Único}}
{
	Arma Mythril & Possui um espaço de Matéria extra. \ofrow
	Fomalhaut & O dano causado é do tipo mágico. \ofrow
	Abelhinha & Num Acerto Crítico, aja novamente. \ofrow
	Metralhadora & Se o alvo se evadir do seu ataque, ele ainda sofre metade do dano. \ofrow
	Baforada de\newline Dragão & O dano causado é do tipo Fogo. Além disso, todos que terminarem o turno dentro de 1u de você, sofrem 3 de dano de Fogo.\ofrow
	Canhão \newline Espiritual & Ao acertar um ataque em um alvo que esteja Zumbi, será um Acerto Crítico.\ofrow
	\ofequipnote{Armas de Fogo tem alcance 3u. Diferentemente dos Arcos, seus danos não são aumentados pela FOR, ao invés disso, o dano delas é aumentado em 1d para cada 3 FOR do usuário.}
}
%
\vfill
%
\oftable{p{0.22\columnwidth} p{0.7\columnwidth}}
{\oficonweapon\accf{Lança} & \accf{Efeito Único}}
{
	Lança\newline Mythril & Possui um espaço de Matéria extra. \ofrow
	Gae Bolg & Tenha um Acerto Crítico quando o alvo rolar 3 ou menos no teste de evasão.\ofrow
	Vel & O dano causado por suas Técnicas se torna do tipo mágico. \ofrow
	Longinus & Quando o alvo do seu ataque rolar 5 ou menos no teste de evasão, ele fica Mudo por 3 rodadas.\ofrow 
	Tridente & Ataca o alvo e qualquer um atrás dele. \ofrow
	Naginata & Num Acerto Crítico, triplique, ao invés de dobrar, o dano causado. \ofrow
	\ofequipnote{Lanças têm 2u de alcance.}
}
%
\newpage
%
\oftable{p{0.22\columnwidth} p{0.7\columnwidth}}
{\oficonweapon\accf{Adaga} & \accf{Efeito Único}}
{	
	Faca Mythril & Possui um espaço de Matéria extra. \ofrow
	Sombra Noturna & Se o alvo rolar 5 ou menos no teste de evasão, ele fica Envenenado por 3 rodadas.\ofrow
	Faca\newline Assassina & Num Acerto Crítico, o alvo sofre KO. \ofrow
	Oricalco & Se seu PV cair a metade do máximo, ganhe AuFOR por 3 rodadas. \ofrow
	Gladio  & A DF de todos os testes relacionados a roubar tem -1. \ofrow
	Adaga de\newline Aparar & DEF~+1. Fornece DEF~+1 extra se equipado como faca secundária. \ofrow
	Sorriso Sutil & Se houver um aliado a 1u do alvo, causa +3 de dano com essa arma. \ofrow
	\ofequipnote{Você pode equipar uma Adaga secundária no lugar de um Acessório. Neste caso, você pode atacar com a arma secundária ao invés de se mover. \phantom{y}}
}
%
\vfill
%
\oftable{p{0.22\columnwidth} p{0.7\columnwidth}}
{\oficonweapon\accf{Espada} & \accf{Efeito Único}}
{
	Espada Mythril & Possui um espaço de Matéria extra. \ofrow
	Espada Bruta & Se o alvo rolar 5 ou menos no teste de evasão, o dano causado ignora a DEF. \ofrow
	Lâmina Pistola & Você pode fazer um ataque à distancia de até 3u. Neste caso, não adicione sua FOR ao dano.\ofrow
	Organyx & Ao acertar um ataque, recupere 2 PM. \ofrow
	Espada\newline Sanguinária & Ao reduzir o alvo a 0 PV, recupere PV igual ao seu nível atual.\ofrow
	Murasame & O alcance de suas Técnicas é aumentado em 1u. \ofrow
	Yoshihara & Ao fazer um Contra-ataque através de um teste de evasão, sempre será Acerto Crítico.\ofrow
	Espada Vorpal & Num Acerto Crítico, triplique, ao invés de dobrar, o dano causado. \ofrow
	Salve a \newline Rainha & Sempre que você ou um aliado até 1u for afetado por uma magia, você pode reduzir o dano pela metade ao passar num teste de DF~9. \ofrow
	\ofequipnote{Ao fazer o teste de evasão, você pode contra-atacar no 11 e 12.}
}
%
\clearpage
%
\oftable{p{0.22\columnwidth} p{0.7\columnwidth}}
{\oficonweapon\accf{Cajado} & \accf{Efeito Único}}
{
	Cajado\newline Mythril & Possui um espaço de Matéria extra. \ofrow
	Cajado\newline Curativo & Ao restaurar PV com magia, adicione +2. \ofrow 
	Cajado de\newline Carvalho & Após acertar um ataque, você pode usar a ação Defender por 4 PM. \ofrow
	Cajado Rúnico & Ao causar dano mágico ao inimigo que sofra de um Estado, ele também fica Mudo por 3 rodadas. \ofrow
	Cajado\newline Dourado & Ao ter seus PV abaixo da metade, ganhe AuDEF por 3 rodadas. \ofrow
	Cajado\newline Poderoso & Após acertar um ataque, adicione também sua MAG ao dano. \ofrow
	Cajado do\newline Sábio & Você pode adicionar o tipo Sagrado aos danos causados por ataques e magias. \ofrow
	Osso Oráculo & O alcance de todas as magias de cura ou Estados aumenta em 1u.\ofrow
	\ofequipnote{Cajados aumentam o PM máximo do usuário em 10 por nível de equipamento. Em troca, eles não podem causar mais do que 1d de dano.}
}
%
\vfill
%
\oftable{p{0.22\columnwidth} p{0.7\columnwidth}}
{\oficonweapon\accf{Cetro} & \accf{Efeito Único}}
{
	Cetro Mythril & Possui um espaço de Matéria extra. \ofrow
	Cetro\newline Elemental & Tipo elemento (Ex.: Cetro de Fogo). Ao causar dano do seu elemento, adicione +2 ao dano. \ofrow
	Cetro\newline Adamante & O dano da arma ignora a DEF do alvo. \ofrow
	Cetro Estelar & Ao reduzir o inimigo a 0 PV, recupere PM igual ao seu nível atual.\ofrow
	Cetro\newline Malboro & Ao causar um ou mais Estados negativos, recupere 3 PV. \ofrow
	Cetro Lilith & Ao ter seu PV abaixo da metade, ganhe AuMAG por 3 rodadas. \ofrow
	Bigode de\newline Baleia & Você pode adicionar o tipo Água aos danos causados por ataques e magias.\ofrow
	Cetro Magus & Ao se beneficiar de AuMAG, o efeito é dobrado. \ofrow 
	\ofequipnote{Cetros aumentam a MAG do usuário em +2 por nível de equipamento, mas não podem causar mais do que 1d de dano.}
}
%
%
\newpage
%
\oftable{p{0.35\columnwidth} p{0.62\columnwidth}}
{\oficonarmor\accf{Armadura\newline Pesada} & \accf{Efeito Único}}
{
	Armadura Mythril & Possui um espaço de Matéria extra.  \ofrow
	Malha de Cristal & Resistência: Gelo\ofrow
	Malha Demoníaca & PM Máximo +5  \ofrow
	Armadura\newline Diamantina & Resistência: Raio  \ofrow
	Malha Dracônica & Resistência: Fogo \ofrow
	Armadura de\newline Cavaleiro & RES +1 \ofrow
	Malha Espelhada & Imunidade: Mudo \ofrow
	Aquiles & Imunidade: Lento \ofrow
	Armadura Completa & PV Máximo +5 \ofrow
	Karuta-gane & Ao Correr, mova-se o dobro. \ofrow
	\ofequipnote{Armaduras Pesadas fornecem DEF~+2 extra.}
}
%
\vfill
%
\oftable{p{0.35\columnwidth} p{0.62\columnwidth}}
{\oficonarmor\accf{Armadura Leve} & \accf{Efeito Único}}
{
	Colete Mythril & Possui um espeço de Matéria extra. \ofrow
	Vestes Gaia & Resistência: Terra \ofrow
	Kenpo Gi & Imunidade: Cego \ofrow
	Minerva & PM Máximo +5 \ofrow
	Colete Miragem & Imunidade: Adormecido \ofrow 
	Vestes Ninja & Imunidade: Imóvel \ofrow			 
	Colete Poderoso & DEF +1 \ofrow
	Jaqueta Rubra & Resistência: Fogo \ofrow
	Colete de\newline Sobrevivência &PV Máximo +5 \ofrow
	Colete à prova\newline de balas & Resistência contra danos de Arcos e Armas de Fogo. \ofrow
	Pele do leão \newline da Neméia & Resistência contra danos de Espadas e Adagas. \ofrow
	Vestes Behemoth & RES +1 \ofrow
	\ofequipnote{Armaduras Leves fornecem DEF~+1 e RES~+1 extras.\phantom{y}}
}
%
\vfill
%
\oftable{p{0.35\columnwidth} p{0.62\columnwidth}}
{\oficonarmor\accf{Robes} & \accf{Efeito Único}}
{
	Robe Mythril &  Possui um espaço de Matéria extra. \ofrow
	Robe Negros & Imunidade: Envenenado \ofrow
	Robe de Seda & DEF +1 \ofrow
	Robe de Algodão & Resistência: Vento \ofrow
	Robe Luminoso & Resistência: Sagrado \ofrow
	Robe Magus & RES +1 \ofrow
	Robe do Acadêmico  & PM Máximo +5  \ofrow
	Mortalha\newline Necromântica & Resistência: Escuridão \newline \ofrow
	Robe Branco & Imunidade: Adormecido \ofrow
	Hagoromo & Imunidades a Campos de Efeito. \ofrow
	\ofequipnote{Robes fornecem RES~+2 extra.}	
}
%
\clearpage
%
\oftablewide{p{0.32\columnwidth} p{0.25\columnwidth} p{1.45\columnwidth}}
{\oficonaccessory\accf{Acessórios} & \accf{Nível} & \accf{Efeito Único}}
{
	 Escudo Mythril & Iniciante & DEF +1  \ofrow
	 Bracelete Poderoso & Iniciante & FOR +1 \ofrow
	 Braçadeiras\newline Rúnicas & Iniciante & RES +1 \ofrow
	 Anel de Cristal & Iniciante & MAG +1 \ofrow
	 Botas de batalha & Iniciante & Imunidade: Imóvel  \ofrow
	 Óculos Prata & Iniciante & Imunidade: Cego  \ofrow
	 Colar Estelar & Iniciante & Imunidade: Envenenado  \ofrow
	 Capa Branca & Iniciante & Imunidade: Mudo  \ofrow
	 Insignia da Fúria & Iniciante & Ganhe AuFOR enquanto seu PM for 5 ou menos. \ofrow
	 Anel de Força & Iniciante & Ganhe AuMAG enquanto seu PM estiver cheio. \ofrow
	 Ninja Tabi & Iniciante & Ao Correr, mova-se o dobro. \ofrow
	 Colar Vital & Iniciante & Ao se recuperar de KO, fique Regenerando por 3 rodadas. \ofrow
	 Cinto Prático & Iniciante & Você pode Re-Equipar como ação livre.\ofrow
	 Broquel & Iniciante & Ao usar Defesa, recupere 3 PV. \ofrow
	 Botas Furtivas & Avançado & Seus passos são inaudíveis. \ofrow
	 Luvas de Ladrão & Avançado & Vantagem em testes relacionados a roubar. \ofrow
	 Anel Protetor & Avançado & Se um ataque te acertar, ganhe AuDEF por 1 rodada. \ofrow
	 Abotoadura \newline Elem. & Avançado & Tipo elemental (ex.: Abotoadura de Fogo). Ao causar danos desse elemento, cause +3 extra.  \ofrow
	 Guarda item & Avançado & Se fora de combate, coloque um item nele. Durante, use o item dele livremente.\ofrow
	 Diadema & Avançado & RES +1, MAG +1\ofrow
	 Elmo Grande & Avançado & FOR +1, DEF +1\ofrow
	 Bandana & Avançado & RES~+1, Imunidade: KO \ofrow 
	 Cinto Campeão & Avançado & FOR +1, Imunidade: ReATR \ofrow
	 Botas Germinas & Avançado & A altura de seus pulos é dobrada. \ofrow
	 Faixa Preta & Avançado & PV Máximo +10  \ofrow
     Anel Coração & Avançado & PM Máximo  +10  \ofrow
	 Amuleto Moogle & Avançado & Brilha quando há um monstro a 50u de você. \ofrow
	 Cinto Musculoso & Avançado & Se o alvo de seu ataque rolar 3 ou menos no teste de evasão, será um Acerto Crítico.\ofrow
 	 Escudo Égide & Especialista & Ao ser alvo de uma Magia que não envolva um teste, você pode fazer um teste de DF~10 e evitar o efeito. \ofrow
	 Pedra Roseta & Especialista & Você é capaz de entender qualquer língua escrita e falada. \ofrow
	 Sapato Alado & Especialista & Em toda batalha, você fica Acelerado no seu primeiro turno. Imunidade: Lento.\ofrow
	 Escudo Heroico & Especialista & DEF +1, RES +1, Imunidade: Adormecido. \ofrow
	 Pulseira Protetiva & Especialista & Ganhe AuDEF e AuRES enquanto estiver com PV abaixo de 50\%. \ofrow
	 Botas Pena & Especialista & Você pode levitar até 1u acima do chão. \ofrow
	 Sandália Alada& Especialista & AGI~+1 \ofrow
	 Sandália Divina & Especialista & Permite andar sobre a água ou líquidos. \ofrow
	 Elmo Genji & Especialista & RES~+2, DEF~+1.\ofrow
	 Escudo Genji & Especialista & DEF~+2, RES~+1.\ofrow
	 Luvas Genji & Especialista & FOR~+2, MAG~+2.\ofrow
	 Grampo Dourado & Especialista & Habilidades tem custo em PM -2.  \ofrow
 	 Fita de cabelo & Especialista & Ganhe Vantagem em todos os testes para resistir a Efeitos de Estado.\ofrow
	 \multicolumn{3}{p{2\columnwidth}}{\hspace*{-0.15cm}\accf{Nota:} Para alguns acessórios não faz sentido usar dois ao mesmo tempo (Ex.: escudos)}
}
%
\clearpage
%
\oftablewide{p{0.3\columnwidth} p{0.18\columnwidth} p{1.5\columnwidth}}
{\oficonitem\accf{Item} & \accf{Valor} & \accf{Efeito Único}}
{
	Verdura Gysahl & 25G & Vegetais favoritos das Chocobos. \ofrow
	Antídoto & 50G & Remove Envenenado. \ofrow
	Colírio & 50G & Remove Cego.  \ofrow 
	Grama Eco & 50G & Remove Mudo.  \ofrow 
	Agulha dourada & 50G & Remove Imóvel. \ofrow
	Vento Ártico & 100G & O alvo sofre 2d da dano de Gelo. \ofrow
	Granada & 100G & O alvo sofre 2d da dano de Fogo. \ofrow
	Gema Raio & 100G & O alvo sofre 2d da dano de Raio. \ofrow
	Poção & 100G & O alvo recupera 8 PV. \ofrow
	Água Benta & 150G & Remove Zumbi.\ofrow
	Éter & 150G & O alvo recupera 12 PM. \ofrow
	Cortina de luz & 200G & O alvo ganha AuDEF por 3 rodadas. \ofrow
	Cortina Lunar& 200G & O alvo ganha AuRES por 3 rodadas. \ofrow
	Tônico Gigante & 200G & O alvo ganha AuFOR por 3 rodadas.\ofrow
	Tônico Feérico & 200G & O alvo ganha AuMAG por 3 rodadas.\ofrow
	Óleo Elemental & 250G & Tipo elemental (Ex.: Óleo Flamejante). Aplique-o na sua arma e seu próximo ataque causa 6 de dano desse elemento.\ofrow
	Remédio & 250G & Remove todos os Estados negativos, exceto KO. \ofrow
	Vinha Malboro & 250G & O alvo faz um teste de DF~8 ou fica Envenenado por 3 rodadas.\ofrow
	Pó do Sono & 250G & O alvo faz um teste de DF~8 ou fica Adormecido por 3 rodadas. \ofrow
	Bebida Heroica & 300G & O alvo ganha AuFOR e AuMAG por 3 rodadas. \ofrow
	Pedra Teleporte & 300G & Teleporte-se a um lugar que possa ver até 10u. \ofrow
	Teia de Aranha & 350G & O alvo faz um teste de DF~8 ou fica Lento por 3 rodadas. \ofrow
	Poção Maior & 400G & O alvo recupera 20 PV. \ofrow
	Escaneadora & 500G & Revele as Resistências, Fraquezas,Imunidades e PV e PM atuais do alvo. \ofrow
	Super Éter & 500G & O alvo recupera 30 PM. \ofrow
	Vacina & 500G & Remove Envenenado e o alvo fica Imune a isso até o fim da batalha. \ofrow
	Megafone & 500G & Remove Mudo e o alvo fica Imune a isso até o fim da batalha. \ofrow
	Lentes Mágicas & 500G & Remove Cego e o alvo fica Imune a isso até o fim da batalha. \ofrow
	Analgésico & 500G & Remove Imóvel e o alvo fica Imune a isso até o fim da batalha. \ofrow
	Pena da Fênix & 500G & Remove KO e o alvo recupera 1 PV. \ofrow
	Matéria Escura & 500G & O alvo sofre 6d de dano de Escuridão. \ofrow
	Tenda & 750G & Permite ao grupo dormir no ermo confortavelmente. \ofrow
	Super Poção & 750G & O alvo recupera todo o PV. \ofrow
	Transfusão & 800G & Remove Zumbi e o alvo fica Imune a isso até o fim da batalha. \ofrow 
	Ampulheta\newline Prateada & 800G & Remove Lento e o alvo fica Acelerado por 3 rodadas. \ofrow
	Mega Poção & 800G & Todos a 1u recuperam 25 PM. \ofrow
	Super Éter & 900G & O alvo recupera todo o PM. \ofrow
	Lasca de Magicita & 1000G & Role 1d e faça uma invocação: 1-Ifrit, 2-Shiva, 3-Ramuh, 4-Odin, 5-Fênix, 6-Bahamut.\ofrow
	Magicita Esper & 1250G & Invoca um Esper específico (Ex.: Magicita Ifrit). Ao usar, ativa-se a habilidade Invocar. \ofrow
	Elixir & 1250G & O alvo recupera todo o PV e PM. \ofrow
	Ampulheta\newline Dourada & 1500G & Todos até 5u, exceto você, ficam Adormecidos por 1 rodada. \ofrow
	Mega Elixir & 1750G & Todos até 1u recuperam todo o PV e PM.\ofrow
	Mega fênix & 2000G & Remove KO de todos até 3u e recupera completamente todo o PV e PM deles.
}
%
\clearpage
%
\oftablewide{p{0.36\columnwidth} p{0.16\columnwidth} p{1.45\columnwidth}}
{\accf{Matéria de Comando} & \accf{Valor} & \accf{Efeito Único}}
{
	Matéria Conjuradora & 250G & Crie um efeito mágico simples, como um barulho alto ou clarão de luz, com uma ação. \ofrow
	Matéria Fogo & 500G & Permite conjurar a habilidade "Fogo" (veja Mago Negro) \ofrow
	Matéria Gelo & 500G & Permite conjurar a habilidade "Nevasca" (veja Mago Negro) \ofrow
	Matéria Raio & 500G & Permite conjurar a habilidade "Raio" (veja Mago Negro) \ofrow  
	Matéria Médica & 500G & Permite conjurar a habilidade "Primeiros Socorros" (veja Sentinela) \ofrow   
	Matéria Fúria & 500G & Permite conjurar a habilidade "Surrar" (veja Guerreiro)\ofrow   
	Matéria Bravura & 500G & Permite conjurar a habilidade "Escapar" (veja Ladrão). \ofrow    
	Matéria Kung-Fu & 500G & Permite conjurar a habilidade "Chute" (veja Monge). \ofrow 
	Matéria Improviso & 750G & Permite conjurar a habilidade "Improviso" (veja Bardo). \ofrow
	Matéria Miragem & 750G & Permite conjurar a habilidade "Imagem" (veja Invocador). \ofrow
	Matéria Cegueira & 500G & Permite conjurar a habilidade "Cegar" (veja Mago Vermelho). \ofrow  
	Matéria Restauração & 750G & Permite conjurar a habilidade "Cura" (veja Mago Branco). \ofrow
	Matéria Armadilha & 750G & Permite conjurar a habilidade "Armar Armadilha" (veja Atirador). \ofrow
	Matéria Dracônica & 750G & Permite conjurar a habilidade "Baforada de Fogo" (veja Dragoon). \ofrow
	Matéria de Aviso & 1000G & Use sua ação para disparar um clarão de luz no ar que é visível até 1000u.\ofrow 
	Matéria Transfusão & 1000G & Use sua ação para reduzir seu PV igual ao seu Nível atual e recuperar PV do alvo até 3u na mesma proporção.\ofrow 
	Matéria Bomba & 1250G & Use sua ação para empurrar todos até 1u de você até 2u. \ofrow 
	Matéria Vácuo & 1250G & Use sua ação para puxar inimigos dentro de 3u em sua direção até 2u. \ofrow
	Matéria Voadora & 1500G & Permite usar a habilidade "Propelir"  (veja Engenhoqueiro).\ofrow
	Matéria Temporal & 1500G & Permite usar a habilidade "Acelerar"  (veja Mago do Tempo).\ofrow   
	Item Matéria & 1500G & Permite usar a habilidade "Bolsos Rasos"  (veja Ladrão).\ofrow   
	Matéria Ampliadora & 1500G & Conjure lupas que permitem ver tudo em um local até 5000u como se estivesse ao seu lado. \ofrow
	Matéria Salto & 1750G & Use sua ação para saltar am altura ou distância até 10u. \ofrow
	Matéria Analítica & 2000G & Use sua ação para descobrir o PV restante do alvo dentro de 5u de você.\ofrow 
	Matéria Debandada & 2000G & Pode Correr e atacar imediatamente. \ofrow
	Matéria Esfregão & 2000G & Use sua ação para remover um Campo ativo ao seu lado. \ofrow
	Matéria Remendo & 2000G & Use sua ação para remover um Estado de Efeito de você. \ofrow
	Matéria Troca & 2000G & Use sua ação para trocar sua posição com um aliado dentro de 1u de você. \ofrow
	Matéria Piso\newline Molhado & 2000G & Use sua ação para criar um Campo Derrapante centrado num ponto dentro de 3u de você, de 3u e que dure por 3 rodadas. Não pode ser usado se houver outro campo estiver ativo. \ofrow	
	Matéria Arcana & 2000G & Use sua ação para detectar a presença de magia imbuída em objetos dentro de 10u. \ofrow
	Matéria Lava & 2000G & Use sua ação para criar um Campo Ardente centrada num ponto dentro de 3u de você, de 3u e que dure por 3 rodadas. Não pode ser usado se houver outro campo estiver ativo. \ofrow
	Matéria Telepática & 2000G & Use sua ação para enviar uma mensagem telepática a um alvo dentro de 10u de você. \ofrow 
	Matéria Alarme & 2000G & Use sua ação para criar uma esfera invisível de 5u de diâmetro onde você está. Ela fica por 1 hora e ao alguém a entravessar, ela é destruída, emitindo um barulho alto que pode ser ouvido até 500u. Não pode ser usado se houver outra esfera ativa. \ofrow
	Matéria Enredante & 2500G & Se fora de combate, você pode criar uma armadilha invisível de 2u de diâmetro no chão. Se o inimigo pisar nela, ele se torna incapacitado temporariamente e se você engajar em combate com ele terá uma rodada surpresa. \ofrow 
	Matéria Camaleão & 3000G & Fora de combate, Misture-se ao ambiente. Inimigos só podem te notar se passarem num teste de DF~9. \ofrow 
}
%
\oftablewide{p{0.36\columnwidth} p{0.16\columnwidth} p{1.45\columnwidth}}
{\accf{Matéria de Suporte} & \accf{Valor} & \accf{Efeito Único}}
{
	Matéria Aquática & 250G & Permite respirar na água. \ofrow 
	Matéria Clima & 250G & Equipamento brilha para indicar tempestade ou chuva chegando. \ofrow 
	ATR+ & 500G & Aumente um atributo em +1: FOR, DEF, MAG ou RES.\ofrow 
	Matéria \newline Resistência & 500G &  Tipo elemental (ex.: Fogo). Ganhe resistência ao elemento.\ofrow
	Matéria \newline Elemental & 500G & Tipo elemental (ex.: Fogo). Seus ataques causam dano desse elemento. \ofrow
	Matéria \newline Imunidade & 700G &  Tipo Efeito de Estado (ex.: Cego). Ganhe imunidade a ele. \ofrow
	Matéria Brilhante & 750G & O equipamento brilha, iluminando área de 5u ao seu redor. \ofrow
	Matéria Estado & 750G & Tipo Efeito de Estado (ex.:Envenenado). Ao acertar um ataque e se o alvo rolar 5 ou menos no teste de evasão, ele sofre disso por 3 rodadas.\ofrow
	Matéria\newline Consumidora & 1000G & Ao usar um item em si mesmo, recupere 3 PV.  \ofrow 
	Matéria Chocobo & 1000G & Plane até o chão ao cair de qualquer altura. \ofrow 
	PV+ & 1250G & PV máximo aumenta em 10. \ofrow
	PM+ & 1250G & PM máximo aumenta em 10. \ofrow 
	Matéria Alerta & 1500G & Permite se evadir de ataques enquanto concentrando. \ofrow
	Matéria Moogle & 1500G & A arma ou armadura brilha na presença de monstros até 25u. \ofrow
	Matéria Visão\newline Noturna & 1500G & Tenha visão no escuro.  \ofrow
	Matéria Sentidos & 1750G & Veja o PV restante de todos os inimigos a 1u de você. \ofrow
	Matéria Fúria & 1750G &  Se seu PV estiver abaixo da metade, adicione +3 ao dano de seus ataques.\ofrow
	Matéria Item & 2000G & O alcance de uso de seus itens ganha +1u. \ofrow
	Matéria Robusta & 2000G & Ao resistir a um Estado, o teste tem DF-1. \ofrow
	Matéria \newline Contra-magia & 2000G & Ao sofrer dano de magia que conheça ou outra do mesmo elemento, você pode conjurá-la no atacante de imediato. \ofrow
	Matéria \newline Espectadora & 2000G & Todo dano de área de efeito sofrido, se não for o alvo principal, é reduzido pela metade. \ofrow
	Matéria \newline Concentração & 2250G & Se estiver concentrando e não puder atacar no seu turno, recupere PM igual ao seu nível atual. \ofrow
	Matéria Flutuar & 2250G & Não seja afetado por Campos. \ofrow
	Matéria Raio-X & 2500G & Veja através de paredes de até 1u de espessura se estiver em frente a elas. \ofrow
	Matéria Escalar & 2500G & Ande em superfícies verticais. \ofrow
	Matéria Atrair & 2500G &  Cada monstro dentro de 2u terá você como alvo sempre que possível. Alguns inimigos são imunes a esse efeito, o MJ decide. \ofrow
	Matéria Aprimorar & 2500G & Ao encaixar essa matéria numa arma ou armadura, ela é melhorada para o próximo nível de equipamento. Essa matéria é destruída no processo. \ofrow
	Matéria Revidar & 2500G & Ao rolar 10 ou mais no teste de evasão, você pode atacar o agressor. \ofrow
	Matéria Drenar & 2750G & Ao acertar um ataque, recupere 3 PV. \ofrow
	Matéria Osmose & 2750G & Ao acertar um ataque, recupere 3 PM. \ofrow
	Matéria PV/PM & 2750G &  Seu PV e PM máximos são trocados. \ofrow
	Matéria Rebobinar & 3000G & Ao rolar 2 ou 3 no teste de evasão, role mais uma vez. \ofrow
	Matéria Alcance & 3000G & Aumente o alcance da arma em +1u. \ofrow
	Matéria Ágil & 3000G & Mova-se 1u a mais por turno. \ofrow
	Matéria Fantasma & 3000G & Caminha através das pessoas, monstros ou objetos sólidos. \ofrow
}
%
\clearpage